%==============================================================================
%  Manual de Desarrollador / SDK
%  Gestor de Identidades de Casa Monarca
%
%  Formato segun rubrica (DocuDesarrolloMA2006B2026):
%   - Clase: report
%   - Fuente: 10 pt, Arial (sustituida por Helvetica via paquete helvet)
%   - Hoja: Carta (letterpaper)
%   - Margenes: superior/inferior 20 mm, izquierdo/derecho 25 mm
%   - Interlineado: 5.0 pt (interpretado como espaciado comodo; ver \setstretch
%     y \parskip abajo, ajustables si el profesor pide otro valor exacto)
%   - ~60% del documento es texto (descripciones e ideas)
%
%  Compilacion recomendada: pdflatex (2 pasadas para indice y referencias).
%==============================================================================
\documentclass[10pt,letterpaper]{report}

%------------------------------ Idioma y codificacion -------------------------
\usepackage[utf8]{inputenc}
\usepackage[T1]{fontenc}
\usepackage[spanish,es-noshorthands]{babel} % es-noshorthands evita choques con listings
\usepackage[strings]{underscore}            % permite escribir _ en texto sin escaparlo

%------------------------------ Fuente Arial (Helvetica) ----------------------
\usepackage{helvet}
\renewcommand{\familydefault}{\sfdefault}    % usa sans-serif (Helvetica ~ Arial) por defecto

%------------------------------ Geometria de pagina ---------------------------
\usepackage[letterpaper,top=20mm,bottom=20mm,left=25mm,right=25mm]{geometry}

%------------------------------ Interlineado ----------------------------------
\usepackage{setspace}
\setstretch{1.25}                 % interlineado comodo (ajustable)
\setlength{\parskip}{5pt}         % separacion entre parrafos ~5pt
\setlength{\parindent}{0pt}

%------------------------------ Graficos, color, tablas -----------------------
\usepackage{graphicx}
\graphicspath{{imagenes/}{./}}
\usepackage{xcolor}
\usepackage{booktabs}
\usepackage{tabularx}
\usepackage{array}
\usepackage{enumitem}
\setlist{topsep=2pt,itemsep=1pt,parsep=0pt}

%------------------------------ Codigo (listings) -----------------------------
\usepackage{listings}
\definecolor{codebg}{gray}{0.95}
\definecolor{codekw}{RGB}{0,90,160}
\definecolor{codecm}{RGB}{90,120,90}
\definecolor{codestr}{RGB}{160,60,30}
\lstset{
  basicstyle=\ttfamily\footnotesize,
  backgroundcolor=\color{codebg},
  keywordstyle=\color{codekw}\bfseries,
  commentstyle=\color{codecm}\itshape,
  stringstyle=\color{codestr},
  breaklines=true,
  breakatwhitespace=false,
  columns=fullflexible,
  keepspaces=true,
  showstringspaces=false,
  frame=single,
  framesep=4pt,
  xleftmargin=6pt,
  xrightmargin=4pt,
  aboveskip=8pt,
  belowskip=8pt
}
\lstdefinestyle{bash}{language=bash}
\lstdefinestyle{python}{language=Python}
\lstdefinestyle{plain}{}

%------------------------------ Diagramas (TikZ) ------------------------------
\usepackage{tikz}
\usetikzlibrary{arrows.meta,positioning,shapes.geometric,fit,backgrounds}

%------------------------------ Encabezado / pie ------------------------------
\usepackage{fancyhdr}
\pagestyle{fancy}
\fancyhf{}
\fancyhead[L]{\footnotesize Gestor de Identidades de Casa Monarca}
\fancyhead[R]{\footnotesize Manual de Desarrollador}
\fancyfoot[C]{\footnotesize \thepage}
\renewcommand{\headrulewidth}{0.4pt}
\renewcommand{\footrulewidth}{0.0pt}

%------------------------------ Hipervinculos (al final) ----------------------
\usepackage[hidelinks]{hyperref}
\hypersetup{
  pdftitle={Manual de Desarrollador - Gestor de Identidades de Casa Monarca},
  pdfauthor={Equipo de desarrollo}
}

%==============================================================================
%  DATOS DE PORTADA  (rellenar / confirmar)
%==============================================================================
\newcommand{\manualNombre}{Gestor de Identidades de Casa Monarca}
\newcommand{\manualVersion}{1.0}
\newcommand{\escudoPath}{tecnologico-de-monterrey-blue.png}
\newcommand{\fechaLugar}{Monterrey, Nuevo Le\'on. 12 de junio de 2026}
\newcommand{\carrera}{Ingenier\'ia en Ciencia de Datos y Matem\'aticas}
\newcommand{\curso}{Uso de \'algebras modernas para seguridad y criptograf\'ia (Gpo 601)}
\newcommand{\grupo}{Grupo 601}
% --- Estudiantes (nombre completo y matricula) ---
\newcommand{\estudiantes}{%
\renewcommand{\arraystretch}{1.2}
\begin{tabular}{ll}
\toprule
\textbf{Nombre} & \textbf{Matr\'icula} \\
\midrule
Ana Lidia Arteaga Bret\'on & A00838732 \\
Ricardo Ballesteros Amavizca & A01255358 \\
Daniel Garnelo Mart\'inez & A00573086 \\
Daniela Ortiz Blanco & A00840140 \\
Axel Gabriel Rocha P\'erez & A01067854 \\
\'Alvaro Yeudiel Padilla Lobat\'on & A01665850 \\
\bottomrule
\end{tabular}}
\newcommand{\profesores}{%
\begin{tabular}{c}
\toprule
Alberto F. Mart\'inez \\
Ra\'ul G\'omez Mu\~noz \\
Daniel Otero Fadul \\
Pedro Leonel Olaya Trejos \\
\bottomrule
\end{tabular}}

%==============================================================================
\begin{document}
\selectlanguage{spanish}

%------------------------------ PORTADA ---------------------------------------
\begin{titlepage}
\thispagestyle{empty}
\centering

% Escudo del Instituto (se omite limpio si el archivo no existe)
\IfFileExists{\escudoPath}{%
  \includegraphics[height=2.8cm]{\escudoPath}\\[10mm]
}{%
  \fbox{\parbox[c][2.8cm][c]{0.40\textwidth}{\centering \itshape Escudo del Instituto\\ (colocar \texttt{\escudoPath})}}\\[10mm]
}

{\LARGE\bfseries Manual de Desarrollador}\\[2mm]
{\Large\bfseries \manualNombre}\\[2mm]
{\large Versi\'on \manualVersion}\\[12mm]

{\large\bfseries Equipo de trabajo}\\[3mm]
\estudiantes\\[10mm]

{\large\bfseries Profesores}\\[3mm]
\profesores\\[10mm]

{\large\bfseries Socio Formador}\\[3mm]
\begin{tabular}{c}
\toprule
Casa Monarca \\
\bottomrule
\end{tabular}

\vfill

{\large Instituto Tecnol\'ogico y de Estudios Superiores de Monterrey}\\[1mm]
{\large Escuela de Ingenier\'ia y Ciencias}\\[1mm]
{\large \carrera}\\[5mm]
{\large \curso}\\[1mm]
{\normalsize \fechaLugar}
\end{titlepage}

%------------------------------ INDICE ----------------------------------------
\pagenumbering{roman}
\tableofcontents
\clearpage
\pagenumbering{arabic}

%==============================================================================
\chapter{Descripci\'on General}
%==============================================================================

\section{Prop\'osito del sistema}
El \textbf{Gestor de Identidades de Casa Monarca} es una aplicaci\'on web cuyo prop\'osito es
administrar, de forma segura y trazable, las identidades digitales y los privilegios de acceso del
personal de Casa Monarca, as\'i como proteger la informaci\'on sensible de los expedientes de las
personas migrantes y refugiadas que la organizaci\'on atiende. Aunque la aplicaci\'on tambi\'en
digitaliza el levantamiento de expedientes (denominados \texttt{encuestas} en el c\'odigo), su
columna vertebral es un \emph{stack} de gesti\'on de identidad y seguridad: autenticaci\'on
robusta de contrase\~nas, control de acceso basado en roles (RBAC), infraestructura de clave
p\'ublica (PKI X.509), llaves de acceso WebAuthn (\emph{passkeys}), cifrado de campo con rotaci\'on
de llaves y una bit\'acora de auditor\'ia.

En otras palabras, el sistema responde a una necesidad concreta del Socio Formador: una
organizaci\'on con recursos limitados que custodia datos personales altamente delicados ---cuya
filtraci\'on podr\'ia poner en riesgo f\'isico y legal a personas ya vulnerables--- requiere
mecanismos de identidad y proteccion de datos del nivel de una instituci\'on grande, pero operables
y mantenibles con su propia infraestructura. El gestor de identidades provee exactamente esa capa:
decide \emph{qui\'en} puede entrar, \emph{qu\'e} puede hacer seg\'un su rol, \emph{c\'omo} se prueba
criptogr\'aficamente su identidad para las acciones cr\'iticas, y \emph{queda registro} de cada
operaci\'on relevante.

\section{Alcance}

\subsection*{Qu\'e hace}
\begin{itemize}
  \item \textbf{Autenticaci\'on de usuarios} con hashing \textbf{Argon2id}, bloqueo de cuentas por
  intentos fallidos y verificaci\'on de contrase\~nas filtradas contra la base de datos de
  \emph{Have I Been Pwned} (HIBP) mediante k-anonimato.
  \item \textbf{Control de acceso basado en roles} (RBAC) con cuatro roles
  (\texttt{usuario}, \texttt{operativo}, \texttt{coordinador}, \texttt{admin}) y un conjunto de
  permisos CRUD por rol.
  \item \textbf{Identidad reforzada para roles cr\'iticos} (\texttt{admin}, \texttt{coordinador})
  v\'ia certificados \textbf{X.509} con reto--respuesta firmado (RSA PKCS\#1 v1.5 + SHA-256) y
  \textbf{passkeys WebAuthn}, con derivaci\'on autom\'atica de un certificado a partir de la primera
  passkey (arquitectura N:1).
  \item \textbf{Cifrado de campo} (Fernet) de los datos sensibles de cada expediente, con
  \textbf{rotaci\'on de llaves} basada en un \emph{keyring} con huellas SHA-256 y re-cifrado masivo
  en segundo plano.
  \item \textbf{Flujo de expedientes} con una m\'aquina de estados revisada en cascada por los
  distintos niveles operativos.
  \item \textbf{Solicitudes ARCO} (Acceso, Rectificaci\'on, Cancelaci\'on y Oposici\'on) con un
  pipeline de aprobaci\'on multinivel.
  \item \textbf{Auditor\'ia} (tabla \texttt{logs}) y \textbf{respaldos cifrados} de la base de datos.
\end{itemize}

\subsection*{Qu\'e l\'imites tiene}
\begin{itemize}
  \item No expone una API JSON/REST de prop\'osito general: la interacci\'on es por
  formularios HTML renderizados en el servidor (Jinja2). Solo los flujos WebAuthn y algunos paneles
  de administraci\'on responden JSON consumido por \texttt{fetch()}.
  \item No integra directorios corporativos (LDAP/Active Directory) ni inicio de sesi\'on \'unico
  (SSO); las cuentas se administran localmente.
  \item No usa un m\'odulo de seguridad por hardware (HSM): el material criptogr\'afico reside en
  archivos locales protegidos (\texttt{key.key}, \texttt{certs/}, \texttt{keys/}).
  \item No est\'a dise\~nado para escalado horizontal con estado compartido: usa SQLite como un
  \'unico proceso. Esto condiciona el despliegue (ver el cap\'itulo de Despliegue y el de
  Problemas Conocidos).
  \item El segundo factor para roles cr\'iticos es exclusivamente certificado X.509 + passkey; no
  hay 2FA por SMS o correo.
\end{itemize}

\section{Casos de uso principales}
\begin{enumerate}
  \item \textbf{Inicio de sesi\'on seguro.} Un colaborador se autentica con usuario y contrase\~na;
  si su rol es cr\'itico, el sistema le exige completar y usar su certificado/passkey antes de
  operar.
  \item \textbf{Levantamiento y canalizaci\'on de un expediente.} Un \texttt{usuario} crea un
  expediente (estado \texttt{borrador}); este avanza por revisi\'on de \texttt{operativo},
  \texttt{coordinador} y, finalmente, \texttt{admin}, cifrando los datos sensibles en cada paso.
  \item \textbf{Administraci\'on de identidades.} Un \texttt{admin} emite, descarga o revoca
  certificados, gestiona passkeys y consulta el estado del v\'inculo passkey--certificado.
  \item \textbf{Rotaci\'on de llave de cifrado.} Un \texttt{admin}, autorizando la acci\'on con su
  passkey, configura y activa una nueva llave; el sistema re-cifra los expedientes en segundo plano.
  \item \textbf{Solicitud ARCO.} Una persona presenta una solicitud sobre sus datos en un formulario
  p\'ublico; la solicitud escala por niveles hasta su resoluci\'on o reenv\'io a \texttt{admin}.
  \item \textbf{Auditor\'ia.} Un \texttt{admin} revisa la bit\'acora de eventos de seguridad y
  operaci\'on.
\end{enumerate}

%==============================================================================
\chapter{Arquitectura del Sistema}
%==============================================================================

\section{Tipo de arquitectura}
El sistema sigue una arquitectura \textbf{mon\'olito web cliente--servidor} renderizada en el
servidor (\emph{server-side rendering}). Es una aplicaci\'on Flask de \emph{archivo \'unico}: casi
toda la l\'ogica de negocio, el acceso a datos, los \emph{helpers} criptogr\'aficos y las rutas
viven en \texttt{app.py} (alrededor de 5\,600 l\'ineas, unas 40 rutas). No hay una capa de modelos u
ORM separada: el acceso a datos se hace con \texttt{sqlite3} crudo a trav\'es de \texttt{get\_conn()}.
La configuraci\'on por ambiente se centraliza en \texttt{config.py}.

Este estilo se eligi\'o deliberadamente por \textbf{simplicidad operativa}: un solo proceso, una sola
base de datos en archivo, sin dependencias de infraestructura pesada, de modo que el Socio Formador
pueda instalarlo y mantenerlo con recursos modestos.

\section{Componentes principales}
\begin{itemize}
  \item \textbf{Capa web (Flask + Jinja2):} define las rutas \texttt{@app.route}, procesa formularios
  y renderiza plantillas HTML en \texttt{templates/}.
  \item \textbf{Control de acceso (RBAC):} \texttt{PERMISSIONS}, \texttt{has\_permission()},
  \texttt{require\_role()} y el \emph{gating} de identidad \texttt{enforce\_certificate\_setup()}.
  \item \textbf{Subsistema de identidad reforzada:} PKI X.509 (CA local en \texttt{certs/}) y
  WebAuthn/passkeys (\texttt{passkey\_credentials}).
  \item \textbf{Subsistema de cifrado:} Fernet con \emph{keyring} de llaves
  (\texttt{encrypt\_data}/\texttt{decrypt\_data}), m\'etricas y cola de re-cifrado.
  \item \textbf{Persistencia:} base de datos SQLite (\texttt{database.db}) con 11 tablas, creada y
  migrada por \texttt{init\_db()} mediante \texttt{ensure\_column()}.
  \item \textbf{Auditor\'ia:} funci\'on \texttt{log()} que escribe en la tabla \texttt{logs}.
  \item \textbf{Procesos auxiliares (\texttt{tools/}):} respaldo/restauraci\'on cifrados y
  \emph{worker} de re-cifrado que consume la cola \texttt{reencrypt\_jobs}.
\end{itemize}

\section{Diagrama general}
La Figura~\ref{fig:arq} muestra la relaci\'on entre los componentes: el navegador interact\'ua por
HTTP con \texttt{app.py}, que concentra la l\'ogica y persiste en SQLite; los procesos auxiliares
operan fuera del ciclo de petici\'on/respuesta sobre la misma base de datos y el material de llaves.

\begin{figure}[h]
\centering
\begin{tikzpicture}[
  font=\footnotesize,
  box/.style={draw, rounded corners, align=center, inner sep=5pt, minimum height=9mm},
  node distance=10mm and 16mm
]
  \node[box, fill=blue!8] (browser) {Navegador\\(cliente web)};
  \node[box, fill=green!10, below=of browser, text width=6.6cm]
    (app) {\textbf{app.py} (Flask, $\sim$5\,600 l\'ineas, $\sim$40 rutas)\\[2pt]
            RBAC \textbullet\ CSRF \textbullet\ Cifrado Fernet\\
            PKI X.509 \textbullet\ WebAuthn \textbullet\ Auditor\'ia};
  \node[box, fill=orange!12, below=of app] (db) {SQLite \textemdash\ \texttt{database.db} (11 tablas)};
  \node[box, fill=gray!12, right=of app, text width=3.6cm]
    (tools) {\texttt{tools/}\\ \texttt{backup\_db.py}\\ \texttt{reencrypt\_worker.py}};
  \node[box, fill=gray!8, below=of tools, text width=3.6cm]
    (keys) {Material de llaves\\ \texttt{key.key}, \texttt{certs/}, \texttt{keys/}};

  \draw[<->, thick] (browser) -- (app) node[midway, right]{HTTP / HTML};
  \draw[<->, thick] (app) -- (db) node[midway, right]{\texttt{sqlite3}};
  \draw[<->] (app) -- (tools);
  \draw[<->] (tools) -- (db);
  \draw[<->] (tools) -- (keys);
  \draw[<->] (app.south west) to[bend right=20] (keys.west);
\end{tikzpicture}
\caption{Diagrama general de componentes del Gestor de Identidades.}
\label{fig:arq}
\end{figure}

%==============================================================================
\chapter{Stack Tecnol\'ogico}
%==============================================================================

\section{Lenguajes}
El backend est\'a escrito \'integramente en \textbf{Python} (3.10 o superior). Las plantillas usan
\textbf{HTML} con el lenguaje de plantillas \textbf{Jinja2}, y el lado del cliente emplea
\textbf{JavaScript} est\'andar del navegador para las ceremonias WebAuthn
(\texttt{navigator.credentials.create/get}) y algunas llamadas \texttt{fetch()}. Las consultas a la
base de datos se expresan en \textbf{SQL} (dialecto SQLite).

\section{Frameworks}
\begin{itemize}
  \item \textbf{Flask 2.3.3} \textemdash\ framework web: enrutamiento, sesiones, render de plantillas.
  \item \textbf{Werkzeug 2.3.7} \textemdash\ capa WSGI subyacente de Flask.
  \item \textbf{Jinja2 3.1.2} \textemdash\ motor de plantillas HTML.
\end{itemize}

\section{Bibliotecas}
La Tabla~\ref{tab:stack} resume las dependencias declaradas en \texttt{requirements.txt} y su
funci\'on criptogr\'afica o de soporte.

\begin{table}[h]
\centering
\renewcommand{\arraystretch}{1.25}
\begin{tabularx}{\textwidth}{l l X}
\toprule
\textbf{Biblioteca} & \textbf{Versi\'on} & \textbf{Uso} \\
\midrule
\texttt{cryptography} & $\geq$43.0.3,$<$46 & Fernet, X.509, RSA, \emph{padding} PKCS1v15, hashes SHA-256. \\
\texttt{argon2-cffi} & 23.1.0 & Hashing de contrase\~nas Argon2id. \\
\texttt{webauthn} & 2.3.0 & Ceremonias WebAuthn (registro y autenticaci\'on de passkeys). \\
\texttt{requests} & 2.31.0 & Llamadas HTTP salientes auxiliares. \\
\texttt{pytest} / \texttt{pytest-cov} & 7.4.3 / 4.1.0 & Suite de pruebas y cobertura. \\
\bottomrule
\end{tabularx}
\caption{Bibliotecas principales del proyecto.}
\label{tab:stack}
\end{table}

La base de datos es \textbf{SQLite}, incluida en la biblioteca est\'andar de Python (\texttt{sqlite3}),
por lo que no requiere un servidor de base de datos aparte.

%==============================================================================
\chapter{Estructura del Proyecto}
%==============================================================================

\section{Organizaci\'on de carpetas}
\begin{lstlisting}[style=plain]
casa-monarca-app/
|- app.py              # Backend completo: rutas, esquema DB, criptografia, reglas
|- config.py           # DevelopmentConfig / ProductionConfig / TestingConfig
|- generate_key.py     # Genera la llave Fernet maestra (key.key)
|- requirements.txt
|- setup.sh            # Instalacion: .venv, deps, .env, key.key, carpetas
|- .env.example        # Plantilla de variables de entorno
|- .gitignore          # Excluye key.key, keys/, certs/*.pem, *.db, etc.
|- templates/          # 14 plantillas Jinja2 (HTML)
|- static/             # Activos estaticos (CSS/JS/imagenes)
|- tools/
|   |- backup_db.py       # Respaldo cifrado de database.db (Fernet)
|   |- restore_db.py      # Restauracion desde respaldo cifrado
|   |- reencrypt_worker.py# Procesa la cola reencrypt_jobs (rotacion)
|- tests/              # Suite pytest (fixtures con DB/certs/llaves aislados)
|- certs/              # CA local + certificados emitidos (gitignored)
|- keys/               # Material de llaves rotadas (gitignored)
|- backups/            # Respaldos cifrados
|- logs/               # Logs de ejecucion
|- key.key             # Llave Fernet activa (gitignored)
|- database.db         # Base de datos SQLite (gitignored, creada al arrancar)
\end{lstlisting}

\section{Descripci\'on de m\'odulos}
\begin{itemize}
  \item \textbf{\texttt{app.py}} \textemdash\ n\'ucleo de la aplicaci\'on. Para orientarse conviene
  buscar \texttt{@app.route} (rutas) y los \texttt{def} de nivel superior (helpers). Agrupa: esquema y
  migraciones (\texttt{init\_db}, \texttt{ensure\_column}), RBAC, contrase\~nas (Argon2id, pol\'itica,
  HIBP), PKI X.509, WebAuthn/passkeys, cifrado Fernet, CSRF y banderas de cookies, y la l\'ogica de
  expedientes y ARCO.
  \item \textbf{\texttt{config.py}} \textemdash\ jerarqu\'ia \texttt{Config} $\rightarrow$
  \texttt{DevelopmentConfig}/\texttt{ProductionConfig}/\texttt{TestingConfig}, seleccionada con la
  variable \texttt{FLASK\_ENV}.
  \item \textbf{\texttt{generate\_key.py}} \textemdash\ genera la llave maestra Fernet en \texttt{key.key}.
  \item \textbf{\texttt{tools/}} \textemdash\ utilidades fuera del ciclo de petici\'on: respaldo y
  restauraci\'on cifrados y el \emph{worker} de re-cifrado.
  \item \textbf{\texttt{templates/}} \textemdash\ 14 plantillas Jinja2 (login, dashboard, survey,
  admin, admin\_cifrado, admin\_identidades, admin\_pki, arco, cert\_setup, colaborador, logs,
  password\_update, profile, usuarios).
  \item \textbf{\texttt{tests/}} \textemdash\ pruebas \texttt{pytest} con \emph{fixtures} que
  reconstruyen una base de datos, certificados y llaves aislados por prueba.
\end{itemize}

%==============================================================================
\chapter{Configuraci\'on del Entorno}
%==============================================================================

\section{Requisitos (versiones)}
\begin{itemize}
  \item \textbf{Python} 3.10 o superior.
  \item Las dependencias de \texttt{requirements.txt} (ver Tabla~\ref{tab:stack}).
  \item Sistema operativo tipo Unix (macOS o Linux) para \texttt{setup.sh}; en Windows pueden
  ejecutarse los pasos manualmente dentro de un entorno virtual.
\end{itemize}

\section{Variables de entorno}
La configuraci\'on se controla con variables de entorno (plantilla en \texttt{.env.example}). Las m\'as
relevantes, con sus valores por defecto, se listan en la Tabla~\ref{tab:env}.

\begin{table}[h]
\centering
\renewcommand{\arraystretch}{1.2}
\begin{tabularx}{\textwidth}{l l X}
\toprule
\textbf{Variable} & \textbf{Default} & \textbf{Funci\'on} \\
\midrule
\texttt{FLASK\_ENV} & development & Selecciona la clase de configuraci\'on. \\
\texttt{SECRET\_KEY} & \itshape inseguro & Firma de cookies de sesi\'on (debe cambiarse en prod). \\
\texttt{ENCRYPTION\_KEY\_PATH} & key.key & Ruta de la llave Fernet activa. \\
\texttt{ENCRYPTION\_LEGACY\_KEY\_PATHS} & (vac\'io) & Llaves ``legadas'' para descifrar durante rotaci\'on. \\
\texttt{CERT\_CA\_CERT\_PATH} & certs/ca\_cert.pem & Certificado de la CA local. \\
\texttt{CERT\_CA\_KEY\_PATH} & certs/ca\_key.pem & Llave privada de la CA local. \\
\texttt{LOGIN\_MAX\_ATTEMPTS} & 5 & Intentos antes de bloquear. \\
\texttt{LOGIN\_WINDOW\_SECONDS} & 300 & Ventana de conteo de intentos. \\
\texttt{LOGIN\_LOCKOUT\_SECONDS} & 900 & Duraci\'on del bloqueo. \\
\texttt{PASSWORD\_MIN\_LENGTH} & 12 & Longitud m\'inima de contrase\~na. \\
\texttt{SIGNATURE\_CHALLENGE\_TTL\_SECONDS} & 300 & Vigencia del reto de firma. \\
\texttt{PASSKEY\_ENABLED} & 1 & Habilita WebAuthn. \\
\texttt{PASSKEY\_ENFORCE\_CRITICAL} & 0 & Exige passkey a roles cr\'iticos. \\
\texttt{PASSKEY\_RP\_ID} / \texttt{PASSKEY\_ORIGIN} & localhost / http://localhost:5000 & Identidad de la \emph{Relying Party}. \\
\texttt{ENABLE\_SESSION\_COOKIE\_SECURE} & 1 (prod) & Fuerza cookies \texttt{Secure}. \\
\bottomrule
\end{tabularx}
\caption{Variables de entorno principales.}
\label{tab:env}
\end{table}

\section{Archivos de configuraci\'on}
\begin{itemize}
  \item \textbf{\texttt{.env}} \textemdash\ valores del entorno local (copiado de \texttt{.env.example}
  por \texttt{setup.sh}); est\'a en \texttt{.gitignore}.
  \item \textbf{\texttt{config.py}} \textemdash\ define las tres clases de configuraci\'on. En
  \texttt{TestingConfig}, \texttt{LOGIN\_MAX\_ATTEMPTS} sube a 1000 para no interferir con las pruebas;
  \texttt{ProductionConfig} fuerza \texttt{SESSION\_COOKIE\_SAMESITE="Strict"} y desactiva
  \texttt{DEBUG}.
  \item \textbf{\texttt{key.key}}, \textbf{\texttt{certs/*.pem}}, \textbf{\texttt{keys/}} \textemdash\
  material criptogr\'afico generado en tiempo de ejecuci\'on; \textbf{nunca debe versionarse} (todos
  est\'an en \texttt{.gitignore}).
\end{itemize}

%==============================================================================
\chapter{Instalaci\'on y Ejecuci\'on}
%==============================================================================

\section{Pasos para levantar el sistema}
\begin{enumerate}
  \item Clonar el repositorio y entrar al directorio del proyecto.
  \item Ejecutar el \emph{script} de configuraci\'on: \texttt{./setup.sh}. Este crea el entorno
  virtual \texttt{.venv}, instala dependencias, copia \texttt{.env.example} a \texttt{.env}, genera
  \texttt{key.key} y crea las carpetas \texttt{certs/}, \texttt{backups/} y \texttt{logs/}.
  \item Editar \texttt{.env} con los valores apropiados (sobre todo \texttt{SECRET\_KEY} en
  producci\'on).
  \item Ejecutar la aplicaci\'on: \texttt{.venv/bin/python app.py}. En el primer arranque,
  \texttt{init\_db()} crea/migra el esquema, \texttt{create\_default\_accounts()} siembra cuentas por
  defecto y \texttt{bootstrap\_dev\_certificates()} genera la CA y certificados de desarrollo.
  \item Abrir \url{http://127.0.0.1:5000} en el navegador.
\end{enumerate}

\section{Comandos principales}
\begin{lstlisting}[style=bash]
# Configuracion inicial (entorno, deps, .env, key.key, carpetas)
./setup.sh

# Ejecutar la app (sirve en http://127.0.0.1:5000)
.venv/bin/python app.py

# Verificacion de sintaxis
.venv/bin/python -m py_compile app.py

# Suite de pruebas completa
PYTHONPATH=. .venv/bin/pytest -q

# Una prueba puntual
PYTHONPATH=. .venv/bin/pytest tests/test_password_security.py -q

# Generar/regenerar la llave Fernet de cifrado de datos
.venv/bin/python generate_key.py

# Respaldo / restauracion cifrados de la base de datos
.venv/bin/python tools/backup_db.py
.venv/bin/python tools/restore_db.py
\end{lstlisting}

No existe un comando de \emph{lint}/\emph{format} dedicado: la v\'ia de verificaci\'on establecida es
\texttt{py\_compile} m\'as \texttt{pytest}.

%==============================================================================
\chapter{Base de Datos}
%==============================================================================

\section{Modelo de datos}
La persistencia es \textbf{SQLite} (\texttt{database.db}), accedida con \texttt{sqlite3} crudo a
trav\'es de \texttt{get\_conn()} (con \texttt{row\_factory} para filas tipo diccionario). El esquema
se crea y evoluciona en \texttt{init\_db()} usando el patr\'on \texttt{ensure\_column(conn, tabla,
columna, definici\'on)}, que a\~nade columnas de forma \textbf{idempotente} sobre bases de datos
existentes ---evitando migraciones destructivas. Consta de \textbf{11 tablas}.

\section{Tablas principales}
\begin{table}[h]
\centering
\renewcommand{\arraystretch}{1.2}
\begin{tabularx}{\textwidth}{l X}
\toprule
\textbf{Tabla} & \textbf{Contenido} \\
\midrule
\texttt{usuarios} & Cuentas: \texttt{username}, \texttt{password\_hash} (Argon2id), \texttt{role}, \texttt{must\_change\_password}, estado de identidad reforzada. \\
\texttt{encuestas} & Expedientes: \texttt{estado}, \texttt{datos} (\emph{blob} cifrado Fernet), \texttt{encryption\_key\_fingerprint}, metadatos de revisi\'on. \\
\texttt{logs} & Bit\'acora: \texttt{usuario}, \texttt{accion}, \texttt{categoria} (operacion/seguridad), \texttt{detalle}, \texttt{timestamp}, \texttt{ip\_address}. \\
\texttt{encryption\_keys} & Ciclo de vida de llaves: \texttt{key\_fingerprint} (SHA-256), \texttt{state} (activo/legada/retiring). \\
\texttt{encryption\_metrics} & M\'etricas de latencia/uso de cifrado y descifrado. \\
\texttt{reencrypt\_jobs} & Cola de re-cifrado masivo durante una rotaci\'on de llave. \\
\texttt{solicitudes\_eliminacion} & Solicitudes de eliminaci\'on de expedientes y su resoluci\'on. \\
\texttt{login\_lockouts} & Bloqueo de cuentas: \texttt{intentos}, ventana, \texttt{locked\_until}. \\
\texttt{certificados} & Certificados X.509: \texttt{serial\_number}, \texttt{certificate\_pem}, \texttt{estado}, \texttt{expires\_at}, \texttt{algorithm}. \\
\texttt{passkey\_credentials} & Credenciales WebAuthn: \texttt{credential\_id}, \texttt{public\_key} (COSE), \texttt{sign\_count}, \texttt{cert\_id}. \\
\texttt{solicitudes\_arco} & Solicitudes ARCO con aprobaci\'on por nivel y \texttt{nivel\_actual}. \\
\bottomrule
\end{tabularx}
\caption{Tablas de la base de datos.}
\label{tab:tablas}
\end{table}

\section{Relaciones}
Las relaciones se modelan de forma ligera mediante claves for\'aneas l\'ogicas (SQLite no exige su
declaraci\'on, pero el c\'odigo las respeta):
\begin{itemize}
  \item \texttt{encuestas.usuario\_id} $\rightarrow$ \texttt{usuarios.id} (autor del expediente; uno a
  muchos: un usuario crea muchos expedientes).
  \item \texttt{certificados.user\_id} $\rightarrow$ \texttt{usuarios.id} (uno a muchos hist\'orico,
  pero solo un certificado \texttt{activo} por usuario).
  \item \texttt{passkey\_credentials.user\_id} $\rightarrow$ \texttt{usuarios.id} y
  \texttt{passkey\_credentials.cert\_id} $\rightarrow$ \texttt{certificados.id}. Esta \'ultima es la
  relaci\'on \textbf{N:1} clave: varias passkeys de un usuario se enlazan a un mismo certificado
  derivado.
  \item \texttt{logs.usuario} referencia (por nombre) a \texttt{usuarios.username} para fines de
  auditor\'ia.
  \item \texttt{encuestas.encryption\_key\_fingerprint} corresponde a un
  \texttt{encryption\_keys.key\_fingerprint}, indicando con qu\'e llave se cifr\'o el registro.
\end{itemize}

%==============================================================================
\chapter{API / Interfaces}
%==============================================================================

La aplicaci\'on \textbf{no expone una API JSON/REST de prop\'osito general}: cada ruta renderiza una
plantilla Jinja2 o redirige tras procesar un formulario. Las respuestas JSON se limitan a los
\emph{endpoints} WebAuthn (que siguen el protocolo del navegador) y a algunos paneles de
administraci\'on consumidos por \texttt{fetch()}. Por ello, en este cap\'itulo ``API/Interfaces'' se
documentan las \textbf{rutas HTTP} como interfaz del sistema.

\section{Endpoints}
La Tabla~\ref{tab:rutas} resume las rutas por \'area funcional (40 rutas en total).

\begin{table}[h]
\centering
\renewcommand{\arraystretch}{1.15}
\footnotesize
\begin{tabularx}{\textwidth}{l l X}
\toprule
\textbf{Ruta} & \textbf{M\'etodos} & \textbf{Funci\'on} \\
\midrule
\multicolumn{3}{l}{\textit{Autenticaci\'on e identidad}}\\
\texttt{/} & GET, POST & Login (Argon2id, bloqueo, gating de identidad). \\
\texttt{/auth/passkey/register/options|verify} & POST & Registro WebAuthn. \\
\texttt{/auth/passkey/login/options|verify} & POST & Autenticaci\'on WebAuthn. \\
\texttt{/auth/check-passkey-required} & POST & Indica si la cuenta exige passkey. \\
\texttt{/action/passkey/options|verify} & POST & Firma de acciones cr\'iticas con passkey. \\
\texttt{/certificado/setup} & GET, POST & Alta de certificado X.509 (servidor o CSR). \\
\texttt{/password/update} & GET, POST & Cambio de contrase\~na. \\
\texttt{/logout} & GET & Cierre de sesi\'on. \\
\midrule
\multicolumn{3}{l}{\textit{Expedientes}}\\
\texttt{/dashboard} & GET & Panel seg\'un rol. \\
\texttt{/survey} & GET, POST & Alta de expediente (cifra \texttt{datos}). \\
\texttt{/bandeja} & GET & Pendientes de revisi\'on del rol. \\
\texttt{/encuesta/<id>/avanzar} & POST & Avanza el estado del expediente. \\
\texttt{/solicitar-eliminacion/<id>} & POST & Solicita eliminaci\'on. \\
\texttt{/solicitud/<id>/resolver} & POST & Resuelve la solicitud. \\
\midrule
\multicolumn{3}{l}{\textit{Usuarios, PKI y passkeys}}\\
\texttt{/usuarios}, \texttt{/eliminar-usuario/<id>} & GET/POST & Gesti\'on de cuentas. \\
\texttt{/admin}, \texttt{/admin/identidades}, \texttt{/admin/pki} & GET & Paneles de administraci\'on. \\
\texttt{/admin/pki/revoke-passkey|revoke-certificate} & POST & Revocaciones. \\
\texttt{/certificado/<id>/descargar|revocar} & POST & Descarga/revoca certificado. \\
\midrule
\multicolumn{3}{l}{\textit{Cifrado y rotaci\'on}}\\
\texttt{/admin/cifrado}, \texttt{/admin/cifrado/metrics|jobs} & GET/POST & Estado, m\'etricas y trabajos. \\
\texttt{/admin/keys/configure} & POST & Configura llave candidata (firma passkey). \\
\texttt{/admin/keys/<fp>/activate} & POST & Activa llave y encola re-cifrado. \\
\midrule
\multicolumn{3}{l}{\textit{Auditor\'ia, perfil y ARCO}}\\
\texttt{/logs}, \texttt{/logs/clear} & GET/POST & Bit\'acora y limpieza (solo admin). \\
\texttt{/profile} & GET, POST & Perfil del usuario. \\
\texttt{/arco}, \texttt{/arco/solicitud} & GET/POST & Formulario y alta de solicitud. \\
\texttt{/arco/<id>/aprobar|resolver|resolver-coordinador} & POST & Pipeline multinivel. \\
\bottomrule
\end{tabularx}
\caption{Rutas HTTP (interfaz del sistema).}
\label{tab:rutas}
\end{table}

\section{M\'etodos y par\'ametros}
Las rutas siguen la convenci\'on HTTP est\'andar: \texttt{GET} para vistas y \texttt{POST} para
operaciones que modifican estado. Los par\'ametros se reciben como \textbf{campos de formulario}
(\texttt{request.form}), \textbf{par\'ametros de ruta} (p.\,ej. \texttt{<int:encuesta\_id>}) o, en los
\emph{endpoints} WebAuthn, como \textbf{cuerpo JSON}. Toda solicitud que modifica estado debe incluir
el \textbf{token CSRF} de la sesi\'on (ver el cap\'itulo de Manejo de Errores y Seguridad).

\section{Ejemplos de uso}
\textbf{Avanzar un expediente} (formulario HTML que dispara \texttt{POST /encuesta/<id>/avanzar}):
\begin{lstlisting}[style=plain]
<form method="POST" action="/encuesta/42/avanzar">
  <input type="hidden" name="csrf_token" value="{{ csrf_token }}">
  <button type="submit">Avanzar al siguiente estado</button>
</form>
\end{lstlisting}

\textbf{Iniciar el registro de una passkey} (cliente, JSON):
\begin{lstlisting}[style=plain]
const resp = await fetch('/auth/passkey/register/options', {
  method: 'POST',
  headers: {'Content-Type': 'application/json', 'X-CSRF-Token': csrfToken}
});
const options = await resp.json();
const credential = await navigator.credentials.create({ publicKey: options });
// ... se envia el atestado a /auth/passkey/register/verify
\end{lstlisting}

%==============================================================================
\chapter{Flujos y L\'ogica de Negocio}
%==============================================================================

\section{Procesos clave del sistema}

\subsection*{Inicio de sesi\'on y refuerzo de identidad}
El login (\texttt{/}) valida usuario/contrase\~na con \texttt{verify\_password\_argon2id}. Antes de
comprobar la contrase\~na, \texttt{\_is\_locked()} verifica si la cuenta/IP est\'a bloqueada. Tras un
fallo, \texttt{\_record\_failed()} incrementa el contador; al alcanzar el m\'aximo se fija el bloqueo y
se registra en \texttt{logs} con \texttt{categoria="seguridad"}. Para roles cr\'iticos,
\texttt{enforce\_certificate\_setup()} (\emph{hook} \texttt{before\_request}) obliga a completar y usar
certificado/passkey antes de operar.

\subsection*{Ciclo de vida del expediente}
La m\'aquina de estados (\texttt{next\_status\_for\_role()}) gobierna las transiciones, atando cada
una al rol que corresponde revisar ese estado:
\begin{center}
\texttt{borrador} $\rightarrow$ \texttt{en\_revision\_operativa} $\rightarrow$
\texttt{en\_revision\_coordinacion} $\rightarrow$ \texttt{validado\_coordinacion} $\rightarrow$
\texttt{cerrado}
\end{center}

\subsection*{Cifrado y rotaci\'on de llaves}
\texttt{encrypt\_data()} cifra con la llave activa y antepone la huella de la llave; el formato del
\emph{blob} resultante es:
\begin{lstlisting}[style=plain]
b"fp:" + <huella SHA-256 de la llave> + b"\n" + <token Fernet>
\end{lstlisting}
\texttt{decrypt\_data()} lee la huella, ubica la llave en el \emph{keyring}
(\texttt{get\_cipher\_for\_fingerprint}) ---que incluye la activa y las legadas--- y descifra. La
rotaci\'on configura/activa una nueva llave y encola \texttt{reencrypt\_jobs}, que
\texttt{tools/reencrypt\_worker.py} procesa en lotes.

\subsection*{Solicitudes ARCO multinivel}
\begin{center}
\texttt{/arco} $\rightarrow$ usuarios $\rightarrow$ operativos $\rightarrow$ coordinadores
$\rightarrow$ \{resuelto \textbar\ reenviado a admin\}
\end{center}
Cada nivel registra su aprobaci\'on en columnas dedicadas (\texttt{aprobado\_<nivel>}, \texttt{\_at},
\texttt{\_por}).

\section{Reglas importantes}
\begin{itemize}
  \item Una transici\'on de estado solo procede si el expediente est\'a en el estado que el rol actual
  puede mover; de lo contrario se rechaza.
  \item Las acciones cr\'iticas (rotar llave, revocar identidades) exigen firma con passkey v\'ia
  \texttt{check\_and\_consume\_passkey\_action} (vigencia de 60 s).
  \item Los retos de firma (certificado) son de un solo uso y caducan a los 300 s.
  \item Si un certificado expira, \texttt{check\_certificate\_expiration} desactiva autom\'aticamente
  las passkeys vinculadas (relaci\'on N:1).
  \item Toda contrase\~na nueva pasa por \texttt{validate\_new\_password\_policy} (longitud, entrop\'ia
  y verificaci\'on HIBP).
\end{itemize}

\section{Casos cr\'iticos}
\begin{itemize}
  \item \textbf{Cuenta bloqueada:} tras 5 intentos fallidos en 5 minutos, la cuenta queda bloqueada 15
  minutos; el contador se reinicia al expirar la ventana.
  \item \textbf{Llave activa perdida sin rotaci\'on ordenada:} sustituir \texttt{key.key} sin el flujo
  de rotaci\'on deja ilegibles los datos cifrados con la llave anterior.
  \item \textbf{Certificado comprometido:} debe \emph{revocarse} de inmediato (no basta esperar a su
  expiraci\'on), lo que adem\'as desactiva sus passkeys.
\end{itemize}

%==============================================================================
\chapter{Ejemplos de Uso}
%==============================================================================

\section{Casos t\'ipicos de uso}
\textbf{Crear y canalizar un expediente.} Un \texttt{usuario} entra a \texttt{/survey}, llena el
formulario y lo guarda (estado \texttt{borrador}, datos cifrados). El \texttt{operativo} lo ve en su
\texttt{/bandeja} y lo avanza; luego el \texttt{coordinador} lo revisa y valida; finalmente el
\texttt{admin} lo cierra. Cada paso queda en la bit\'acora.

\section{Integraci\'on b\'asica}
Aunque no hay SDK/cliente externo, los \emph{helpers} criptogr\'aficos pueden usarse de forma
program\'atica dentro del proyecto. Por ejemplo, para cifrar y descifrar un dato sensible:
\begin{lstlisting}[style=python]
from app import encrypt_data, decrypt_data

blob = encrypt_data("Informacion sensible del expediente")
# blob == b"fp:<huella>\n<token Fernet>"

texto = decrypt_data(blob)   # usa el keyring (llave activa + legadas)
assert texto == "Informacion sensible del expediente"
\end{lstlisting}

\section{Ejecuci\'on de funciones clave}
\textbf{Hashing y verificaci\'on de contrase\~na (Argon2id):}
\begin{lstlisting}[style=python]
from app import hash_password_argon2id, verify_password_argon2id

hashed = hash_password_argon2id("UnaContrasenaFuerte!2026")
ok = verify_password_argon2id("UnaContrasenaFuerte!2026", hashed)  # True
\end{lstlisting}

\textbf{Generaci\'on de la llave Fernet} (desde la l\'inea de comandos):
\begin{lstlisting}[style=bash]
.venv/bin/python generate_key.py   # escribe key.key
\end{lstlisting}

%==============================================================================
\chapter{Manejo de Errores}
%==============================================================================

\section{Estrategia de errores}
El sistema no usa un \emph{framework} de errores estructurado tipo JSON. Los errores de negocio se
comunican con \textbf{mensajes \emph{flash}} renderizados en las plantillas
(\texttt{flash()} + \texttt{render\_template}); los flujos WebAuthn/AJAX responden \textbf{JSON} con un
c\'odigo HTTP apropiado y un campo de error legible. Los fallos criptogr\'aficos (firma inv\'alida,
certificado expirado, reto vencido, descifrado fallido) producen un mensaje espec\'ifico al usuario
y, en paralelo, una entrada de auditor\'ia con \texttt{categoria="seguridad"}.

La protecci\'on \textbf{CSRF} es global (\emph{hook} \texttt{before\_request}): \texttt{ensure\_csrf\_token()}
genera el token, \texttt{validate\_csrf()} lo compara y \texttt{csrf\_protect()} rechaza toda
solicitud que modifique estado sin un token v\'alido.

\section{Mensajes comunes}
\begin{itemize}
  \item ``Credenciales incorrectas.'' \textemdash\ usuario o contrase\~na inv\'alidos.
  \item ``Cuenta bloqueada temporalmente.'' \textemdash\ se super\'o \texttt{LOGIN\_MAX\_ATTEMPTS}.
  \item ``Debe configurar su certificado/passkey.'' \textemdash\ \emph{gating} de identidad para rol
  cr\'itico.
  \item ``Firma inv\'alida'' / ``Reto expirado.'' \textemdash\ fallo en reto--respuesta.
  \item ``Contrase\~na insegura o presente en filtraciones.'' \textemdash\ rechazo de pol\'itica/HIBP.
  \item ``Token CSRF inv\'alido.'' \textemdash\ formulario sin token o token incorrecto.
\end{itemize}

\section{Registro (logs)}
La funci\'on \texttt{log(usuario, accion, categoria="operacion", detalle=None)} inserta en la tabla
\texttt{logs} con marca de tiempo e IP. Se invoca en cada operaci\'on sensible: logins
(\'exito/fallo/bloqueo), cambios de estado, emisi\'on/revocaci\'on de certificados y passkeys,
configuraci\'on/rotaci\'on de llaves, resoluciones ARCO y eliminaci\'on de usuarios. La categor\'ia
\texttt{seguridad} separa los eventos de inter\'es forense de las operaciones de negocio
(\texttt{operacion}). La ruta \texttt{/logs} permite visualizar la bit\'acora; \texttt{/logs/clear} su
limpieza (solo \texttt{admin}).

%==============================================================================
\chapter{Pruebas}
%==============================================================================

\section{Tipos de tests}
La suite vive en \texttt{tests/} y usa \emph{fixtures} (\texttt{app\_module}) que aplican
\texttt{monkeypatch.chdir(tmp\_path)} para reconstruir, por cada prueba, una base de datos, un
conjunto de certificados y llaves \textbf{aislados}. La cobertura incluye:
\begin{itemize}
  \item Seguridad de contrase\~nas: rechazo de contrase\~nas d\'ebiles y bloqueo de hashes
  ``legados''.
  \item L\'imite de intentos de login y bloqueo temporal.
  \item Emisi\'on, activaci\'on, revocaci\'on y flujo CSR de certificados X.509.
  \item Firma de acciones cr\'iticas con passkeys y verificaci\'on CSRF.
  \item Ciclo completo de canalizaci\'on de un expediente entre los cuatro roles.
  \item Flujo y escalamiento de solicitudes ARCO multinivel (incluido el reenv\'io a \texttt{admin}).
\end{itemize}
Los archivos \texttt{test\_arco\_*.py}, \texttt{test\_csrf\_arco.py} y \texttt{test\_final\_verification.py}
en la ra\'iz son \emph{scripts} de verificaci\'on ad hoc, no parte de la suite \texttt{pytest} formal.

\section{C\'omo ejecutarlos}
\begin{lstlisting}[style=bash]
# Suite completa
PYTHONPATH=. .venv/bin/pytest -q

# Un archivo
PYTHONPATH=. .venv/bin/pytest tests/test_password_security.py -q

# Una funcion especifica
PYTHONPATH=. .venv/bin/pytest \
  tests/test_password_security.py::test_qa_a_rejects_weak_password_on_forced_update -q
\end{lstlisting}

%==============================================================================
\chapter{Despliegue}
%==============================================================================

\section{C\'omo hacerlo}
\subsection*{Despliegue local / autoalojado (referencia)}
\begin{enumerate}
  \item Ejecutar \texttt{setup.sh} en el servidor de destino.
  \item Ajustar \texttt{.env} para producci\'on: \texttt{FLASK\_ENV=production}, \texttt{SECRET\_KEY}
  fuerte y \'unica, \texttt{ENABLE\_SESSION\_COOKIE\_SECURE=1}.
  \item Ejecutar la app detr\'as de un \emph{proxy} reverso/terminador TLS (la app no incluye servidor
  WSGI de producci\'on ni \emph{proxy} propios).
  \item Resguardar \texttt{key.key}, \texttt{certs/*.pem} y \texttt{database.db} con permisos
  restrictivos.
  \item Programar respaldos cifrados peri\'odicos con \texttt{tools/backup\_db.py} y validar la
  restauraci\'on.
\end{enumerate}

\subsection*{Despliegue en HostGator (hosting de producci\'on)}
En el hosting compartido de \textbf{HostGator}, la aplicaci\'on se publica como aplicaci\'on Python a
trav\'es de \emph{Passenger}/cPanel (``Setup Python App''):
\begin{enumerate}
  \item Crear la aplicaci\'on Python en cPanel, seleccionando la versi\'on de Python (3.10+ si est\'a
  disponible) y el directorio del proyecto.
  \item Instalar dependencias en el entorno virtual que cPanel administra (\texttt{pip install -r
  requirements.txt}).
  \item Definir las variables de entorno desde el panel (las de la Tabla~\ref{tab:env}), en especial
  \texttt{SECRET\_KEY}, \texttt{FLASK\_ENV=production} y los \texttt{PASSKEY\_*} con el dominio real
  (\texttt{PASSKEY\_RP\_ID} y \texttt{PASSKEY\_ORIGIN} deben coincidir con el dominio HTTPS de
  producci\'on, requisito de WebAuthn).
  \item Subir el material de llaves (\texttt{key.key}, CA) por un canal seguro y fuera de la ra\'iz
  p\'ublica.
  \item Apuntar el \emph{passenger\_wsgi} al objeto \texttt{app} de Flask.
\end{enumerate}

\section{Entornos}
Se distinguen tres entornos seleccionados por \texttt{FLASK\_ENV}:
\begin{itemize}
  \item \textbf{development:} \texttt{DEBUG=True}, cookies sin \texttt{Secure}; para trabajo local.
  \item \textbf{production:} \texttt{DEBUG=False}, \texttt{SameSite=Strict}, cookies \texttt{Secure};
  para HostGator.
  \item \textbf{testing:} \emph{rate-limiting} desactivado; para la suite de pruebas.
\end{itemize}

\section{Consideraciones}
\begin{itemize}
  \item \textbf{HTTPS obligatorio para WebAuthn:} las passkeys solo funcionan sobre un origen seguro;
  el dominio de HostGator debe servirse por HTTPS y coincidir con \texttt{PASSKEY\_ORIGIN}.
  \item \textbf{SQLite en hosting compartido:} el bloqueo de archivos puede degradarse bajo
  concurrencia; conviene mantener un solo proceso/worker.
  \item \textbf{Procesos en segundo plano:} el \emph{worker} de re-cifrado
  (\texttt{reencrypt\_worker.py}) puede requerir ejecuci\'on por \emph{cron} si el hosting no permite
  procesos de larga duraci\'on (ver Problemas Conocidos).
  \item \textbf{Secretos fuera del repositorio:} nunca subir \texttt{key.key}, \texttt{certs/*.pem} ni
  \texttt{.env} al control de versiones ni a directorios p\'ublicos.
\end{itemize}

%==============================================================================
\chapter{Seguridad}
%==============================================================================

\section{Autenticaci\'on}
La autenticaci\'on base usa \textbf{Argon2id} (\texttt{Type.ID}) con \texttt{memory\_cost=65536}
(64\,MiB), \texttt{time\_cost=3}, \texttt{parallelism=2}, \texttt{hash\_len=32} y una sal aleatoria de
16 bytes generada con \texttt{os.urandom}. Argon2id es \emph{memory-hard} y resistente tanto a ataques
de canal lateral como a fuerza bruta con GPU/ASIC. Para roles cr\'iticos se a\~nade un segundo factor
de \textbf{posesi\'on}: certificado X.509 con reto--respuesta (RSA, \texttt{PKCS1v15} + \texttt{SHA256})
y/o passkey WebAuthn. Las passkeys son \textbf{resistentes a \emph{phishing}} porque la firma queda
atada al \texttt{PASSKEY\_ORIGIN}, y verifican el \texttt{sign\_count} para detectar clonaci\'on de
autenticadores.

\section{Manejo de datos sensibles}
Los datos sensibles de cada expediente se cifran a nivel de campo con \textbf{Fernet}
(AES-128-CBC + HMAC-SHA256, \emph{cifrado autenticado}). Cada \emph{blob} lleva la \textbf{huella
SHA-256} de la llave usada, lo que permite \textbf{rotaci\'on de llaves} sin perder acceso a datos
antiguos: conviven una llave activa y llaves legadas en el \emph{keyring}. La rotaci\'on sigue el
ciclo activo $\rightarrow$ legada $\rightarrow$ retirada, con re-cifrado masivo en segundo plano. Los
respaldos de la base de datos tambi\'en se cifran con Fernet.

\section{Buenas pr\'acticas}
\begin{itemize}
  \item \textbf{Verificaci\'on de contrase\~nas filtradas} con HIBP por k-anonimato: solo se env\'ian
  los 5 primeros caracteres hex del \texttt{SHA-1} de la contrase\~na; la comparaci\'on del sufijo es
  local.
  \item \textbf{Cookies endurecidas} (\texttt{enforce\_cookie\_flags}): \texttt{HttpOnly},
  \texttt{Secure} y \texttt{SameSite}.
  \item \textbf{CSRF} global con token sincronizado por sesi\'on.
  \item \textbf{M\'inimo privilegio} y \textbf{separaci\'on de funciones} v\'ia RBAC y la m\'aquina de
  estados de expedientes.
  \item \textbf{Secretos fuera del repositorio} (todo el material criptogr\'afico est\'a en
  \texttt{.gitignore}).
  \item \textbf{Auditor\'ia} de eventos de seguridad para no repudio y rendici\'on de cuentas.
\end{itemize}

%==============================================================================
\chapter{Gu\'ia de Contribuci\'on}
%==============================================================================

\section{Convenciones de c\'odigo}
\begin{itemize}
  \item Todo el backend reside en \texttt{app.py}; al agregar l\'ogica, ub\'iquela cerca de las
  funciones afines y respete el estilo existente (nombres en espa\~nol para rutas, columnas y mensajes;
  funciones \emph{snake\_case}).
  \item Para evolucionar el esquema, use \textbf{siempre} \texttt{ensure\_column()} (idempotente) en
  lugar de \texttt{ALTER TABLE} directo o migraciones destructivas.
  \item No imprima ni registre material secreto (llaves, PEM, contrase\~nas).
  \item Verifique con \texttt{py\_compile} y ejecute \texttt{pytest} antes de proponer cambios.
\end{itemize}

\section{Flujo de trabajo (git)}
\begin{enumerate}
  \item Crear una rama a partir de \texttt{main} (no trabajar directamente sobre \texttt{main}).
  \item Realizar \emph{commits} peque\~nos y descriptivos (en espa\~nol, consistente con el historial).
  \item Abrir un \emph{Pull Request} describiendo qu\'e cambia, por qu\'e y c\'omo se prob\'o.
  \item No incluir en el control de versiones \texttt{key.key}, \texttt{certs/*.pem},
  \texttt{keys/}, \texttt{.env} ni \texttt{*.db}.
\end{enumerate}

\section{C\'omo agregar funcionalidades}
\begin{enumerate}
  \item \textbf{Nueva ruta:} a\~nadir un \texttt{@app.route} con el decorador \texttt{require\_role()}
  apropiado y validar permisos con \texttt{has\_permission()}.
  \item \textbf{Nueva tabla/columna:} declararla en \texttt{init\_db()} y, si aplica a bases
  existentes, a\~nadirla con \texttt{ensure\_column()}.
  \item \textbf{Operaci\'on cr\'itica:} protegerla con \texttt{check\_and\_consume\_passkey\_action} y
  registrar el resultado con \texttt{log(..., categoria="seguridad")}.
  \item \textbf{Plantilla:} crear el HTML correspondiente en \texttt{templates/} incluyendo el campo
  \texttt{csrf\_token} en todo formulario \texttt{POST}.
  \item \textbf{Pruebas:} agregar casos en \texttt{tests/} siguiendo el patr\'on de \emph{fixtures}
  aislados.
\end{enumerate}

%==============================================================================
\chapter{Problemas Conocidos}
%==============================================================================

\section{Bugs actuales}
Al tratarse de la primera versi\'on estable del gestor, no hay un registro de \emph{bugs} abiertos
considerados bloqueantes. Las \'areas a vigilar son la concurrencia sobre SQLite y la coordinaci\'on
del \emph{worker} de re-cifrado durante rotaciones de llave.

\section{Limitaciones}
Las principales limitaciones provienen del modelo de despliegue en \textbf{hosting compartido
(HostGator)} y de la arquitectura de archivo \'unico:
\begin{itemize}
  \item \textbf{Soporte de Python/WSGI restringido:} la versi\'on de Python y las extensiones nativas
  disponibles dependen del plan; \texttt{cryptography} y \texttt{argon2-cffi} requieren ruedas
  compatibles. Conviene verificar la versi\'on antes de desplegar.
  \item \textbf{Sin procesos de larga duraci\'on:} el hosting compartido suele matar procesos
  persistentes, por lo que el \emph{worker} de re-cifrado debe ejecutarse de forma programada (cron)
  en lotes, en vez de como demonio continuo.
  \item \textbf{Concurrencia limitada con SQLite:} el bloqueo a nivel de archivo y los l\'imites de
  E/S del hosting compartido restringen el n\'umero de escrituras simult\'aneas; el sistema se piensa
  para un solo proceso.
  \item \textbf{L\'imites de recursos:} cuotas de CPU, memoria y conexiones del plan compartido pueden
  afectar el costo de Argon2id (64\,MiB por verificaci\'on) bajo picos de login.
  \item \textbf{Control parcial del entorno:} la terminaci\'on TLS, las cabeceras y el \emph{cron} los
  administra el panel del proveedor, no la aplicaci\'on; se depende de su configuraci\'on para
  garantizar HTTPS (indispensable para WebAuthn).
  \item \textbf{Sin escalado horizontal:} no hay estado compartido entre instancias; replicar el
  proceso requerir\'ia migrar de SQLite a un motor cliente--servidor.
\end{itemize}

%==============================================================================
\chapter{FAQ T\'ecnica}
%==============================================================================

\section{Problemas comunes al levantar el proyecto}
\begin{description}[leftmargin=0pt, style=nextline]
  \item[\texttt{ModuleNotFoundError} al ejecutar \texttt{app.py}.] Activar el entorno virtual
  (\texttt{source .venv/bin/activate} o usar \texttt{.venv/bin/python}) y reinstalar dependencias con
  \texttt{pip install -r requirements.txt} (o \texttt{setup.sh}).
  \item[La app no encuentra \texttt{key.key}.] Generarla con \texttt{python generate\_key.py} o dejar
  que \texttt{setup.sh} la cree. En HostGator, subirla fuera de la ra\'iz p\'ublica.
  \item[Las passkeys fallan en producci\'on.] Verificar que \texttt{PASSKEY\_RP\_ID} y
  \texttt{PASSKEY\_ORIGIN} coincidan con el dominio HTTPS real; WebAuthn exige origen seguro.
  \item[Cuenta bloqueada en pruebas manuales.] Esperar \texttt{LOGIN\_LOCKOUT\_SECONDS} (15 min) o usar
  el entorno \texttt{testing}, que relaja el \emph{rate-limiting}.
\end{description}

\section{Errores t\'ipicos de desarrollo}
\begin{description}[leftmargin=0pt, style=nextline]
  \item[Datos ilegibles tras regenerar la llave.] Reemplazar \texttt{key.key} sin pasar por el flujo
  de rotaci\'on (\texttt{/admin/keys/configure} + \texttt{/admin/keys/<fp>/activate} +
  \emph{worker}) deja ilegibles los registros previos. Conservar la llave anterior como ``legada''
  hasta completar el re-cifrado.
  \item[``Token CSRF inv\'alido'' en formularios nuevos.] Olvidar incluir
  \texttt{\{\{ csrf\_token \}\}} en un formulario \texttt{POST}.
  \item[Cambios de esquema que rompen bases existentes.] Usar \texttt{ALTER TABLE} directo en lugar de
  \texttt{ensure\_column()}.
  \item[Permisos insuficientes en una ruta nueva.] Olvidar \texttt{require\_role()}/
  \texttt{has\_permission()}, dejando la acci\'on accesible a roles que no deber\'ian ejecutarla.
\end{description}

%==============================================================================
\chapter{Referencias}
%==============================================================================
\begin{itemize}
  \item Documentaci\'on de \texttt{cryptography} (Fernet, X.509, RSA): \url{https://cryptography.io/}
  \item \texttt{argon2-cffi} (Argon2id): \url{https://argon2-cffi.readthedocs.io/}
  \item Biblioteca \texttt{webauthn} (py\_webauthn): \url{https://github.com/duo-labs/py_webauthn}
  \item Especificaci\'on WebAuthn (W3C): \url{https://www.w3.org/TR/webauthn-2/}
  \item Have I Been Pwned \textemdash\ API de contrase\~nas (k-anonimato):
  \url{https://haveibeenpwned.com/API/v3#PwnedPasswords}
  \item RFC 9106 \textemdash\ Argon2: \url{https://www.rfc-editor.org/rfc/rfc9106}
  \item NIST SP 800-57 \textemdash\ Recomendaciones de gesti\'on de llaves:
  \url{https://csrc.nist.gov/publications/detail/sp/800-57-part-1/rev-5/final}
  \item Documentaci\'on de Flask: \url{https://flask.palletsprojects.com/}
  \item Configuraci\'on de aplicaciones Python en HostGator/cPanel (Setup Python App):
  \url{https://www.hostgator.com/help}
\end{itemize}

\end{document}
